% =============================================================================
\section{Population Synthesis}
\label{sec:distribution}
% =============================================================================

A detector finds tens of flowers in a meadow that contains thousands. Placing only
what was detected leaves a scene correct everywhere it was measured and empty
everywhere it was not, which reads worse than either extreme: the eye sees the
sampling pattern rather than the meadow.

Two mechanisms fill that gap, and they are deliberately separate because they
answer different questions. Population synthesis asks \emph{where else would
instances of this class be, given where these ones are}. Ground cover asks
\emph{what is the ground made of}, a question with no discrete instances to
generalize from at all.

\begin{figure}[htbp]
  \centering
  \begin{subfigure}{0.42\linewidth}
    \includegraphics[width=\linewidth]{rainier/synthesis-tree.png}
    \caption{Synthesized tree population.}
    \label{fig:pop-synth}
  \end{subfigure}\hfill
  \begin{subfigure}{0.42\linewidth}
    \includegraphics[width=\linewidth]{rainier/grass-area.png}
    \caption{Ground-cover mask.}
    \label{fig:pop-grass}
  \end{subfigure}
  \caption{The two mechanisms on the Rainier capture, both top-down over the
  \SI{200}{\meter} grid. Population synthesis (a) extends the observed tree
  detections outward across the region where that class is admissible, matching
  their measured spatial statistics rather than repeating a fixed density. Ground
  cover (b) is a mask, not a set of instances: it marks where the scattered tuft
  geometry is admissible. The mask is small relative to the grid because it is
  derived from the region typing of the \emph{original} panorama, which only sees
  meadow in the near field the photograph actually observed.}
  \label{fig:rainier-populations}
\end{figure}

\subsection{Synthesis from Observed Statistics}

For each (class, region-type) pair with enough real detections to constitute
evidence, we measure the spatial statistics of those detections as a
pair-correlation histogram, and synthesize additional instances matching them over
the rest of the matching region. This is a point-process
formulation~\cite{9cfc5dfd-c17c-3458-81d1-9c27d3a321e8}, closely related to
appearance-guided element arrangement~\cite{10.1145/1572614.1572623} and in the
spirit of example-based world synthesis~\cite{10.1145/2766975}.

\subsection{Ground Cover}

Grass is not synthesized from detections, because a detector does not usefully
enumerate blades of grass. It is scattered by density over a mask, and the
interesting part is how that mask is derived.

The mask comes from the region typing of the \emph{original} panorama, for the
reason given in \cref{sec:panorama}: the object-removed panorama has had the
meadow erased along with the objects standing in it. But region type alone is too
blunt in a way that matters here. The semantic label set folds \texttt{grass},
\texttt{earth} and \texttt{field} together with \texttt{sand}, \texttt{snow} and
\texttt{ice} into one ground class, so a snowfield and a meadow are
indistinguishable by type, and an alpine meadow is typically ringed by melting
snowpack that types identically to the wildflowers beside it.

Color settles it. Each candidate cell is tested against the excess-green of the
panorama pixel it was observed through, available directly, because that pixel
is the one retained in \cref{sec:uv-retention}. The threshold is deliberately a low
bar rather than a test for lushness: dry grass sits near zero on this scale while
sand runs about $-0.12$ and snow lower still, so it removes only surfaces that are
definitively not vegetation. The same measurement is then reused as a
\emph{weight} rather than a veto, so density tracks how green the ground actually
is instead of being uniform over whatever survives the threshold.

Tufts are placed as fixed-orientation crossed cards rather than camera-facing
billboards. A billboard is a single quad that turns to face the viewer, which
works for an upright subject at eye level and collapses to a line, then swings
through the ground plane, for ground cover underfoot.
